\documentclass[11pt,a4paper]{amsart}
\usepackage[T1]{fontenc}
\usepackage[utf8]{inputenc}
\usepackage{lmodern}
\usepackage[margin=28mm]{geometry}
\usepackage{amsmath,amssymb,amsthm,mathtools}
\usepackage[expansion=false]{microtype}
\usepackage{enumitem}
\usepackage{xcolor}
\usepackage[colorlinks=true,linkcolor=blue!45!black,
  citecolor=blue!45!black,urlcolor=blue!45!black,
  pdftitle={Simple twisted group algebras over arbitrary fields},
  pdfsubject={Proposed proof of Kourovka Problem 18.78}]{hyperref}

\DeclareMathOperator{\Irr}{Irr}
\DeclareMathOperator{\IBr}{IBr}
\DeclareMathOperator{\Aut}{Aut}
\DeclareMathOperator{\Gal}{Gal}
\DeclareMathOperator{\Rad}{Rad}
\DeclareMathOperator{\GL}{GL}

\newtheorem{theorem}{Theorem}[section]
\newtheorem{proposition}[theorem]{Proposition}
\newtheorem{lemma}[theorem]{Lemma}
\newtheorem{corollary}[theorem]{Corollary}
\theoremstyle{remark}
\newtheorem{remark}[theorem]{Remark}

\numberwithin{equation}{section}
\setlength{\emergencystretch}{2em}
\setcounter{tocdepth}{1}

\title[Simple twisted group algebras]{Simple twisted group algebras\\over arbitrary fields}
\date{26 October 2026}

\begin{document}
\begin{abstract}
A proof is proposed that a finite group admitting a simple twisted group
algebra over an arbitrary field is soluble. The argument establishes a
stronger assertion about transitive actions on irreducible characters
above a faithful central character, allowing group automorphisms and
cyclotomic Galois automorphisms to act together. Ordinary characters
and $2$-Brauer characters are treated by one minimal-counterexample
argument. An affine classification excludes the Brauer case and forces
the ordinary case to have radical quotient a direct product of copies
of $\mathrm{GL}_3(2)$. Canonical lifts give Galois stability in this
class. Restriction to odd-order elements, followed by a central
quotient, then produces a smaller counterexample in characteristic two.
Galois descent and a Witt-vector lift connect the character theorem
to twisted group algebras in every characteristic.
\end{abstract}
\maketitle
\begin{center}
\small\itshape Research preprint. The proposed proof requires independent verification.
\end{center}
\medskip


\section{Introduction}

Let $G$ be a finite group and let $K$ be a field. A normalized
$2$-cocycle is a function $\alpha:G\times G\longrightarrow K^\times$
satisfying
\[
 \alpha(1,g)=\alpha(g,1)=1,\qquad
 \alpha(g,h)\alpha(gh,x)=\alpha(h,x)\alpha(g,hx).
\]
The twisted group algebra $K^\alpha G$ has a $K$-basis
$\{u_g:g\in G\}$ and multiplication
$u_gu_h=\alpha(g,h)u_{gh}$. Simplicity means that its only two-sided
ideals are zero and the whole algebra; the center need not equal $K$.

Passman's question, recorded as Problem~18.78 in the
\emph{Kourovka Notebook} \cite{Kourovka}, asks whether simplicity of
$K^\alpha G$ forces $G$ to be soluble. The distinction between simple
and central simple algebras is essential. In characteristic zero, the
central simple case follows from the theorem on groups of central
type, whereas an arbitrary simple algebra can split into several
simple components after extending the field \cite{Passman}.

The purpose of this article is to give a proposed proof of the following
statement.

\begin{theorem}\label{thm:main}
Let $G$ be a finite group, let $K$ be a field, and let
$\alpha:G\times G\longrightarrow K^\times$ be a normalized $2$-cocycle.
If $K^\alpha G$ is simple, then $G$ is soluble.
\end{theorem}

The character-theoretic statement underlying the proof permits group
automorphisms and Galois automorphisms to act together. Write
$\mathbb Q^{\mathrm{ab}}$ for the union of the cyclotomic fields in
$\mathbb C$, and put
$\Omega=\Gal(\mathbb Q^{\mathrm{ab}}/\mathbb Q)$.
For $a\in\Aut(D)$, $\sigma\in\Omega$, and an ordinary character $\chi$,
the paired action is
\[
 \chi^{(a,\sigma)}(g)=\sigma\bigl(\chi(a^{-1}(g))\bigr).
\]
The two operations commute. If $C\leq Z(D)$ and
$\lambda\in\Irr(C)$ is linear, let
\[
 \Irr(D\mid\lambda)
 =\{\chi\in\Irr(D):\chi|_C=\chi(1)\lambda\}.
\]
Here $\Irr(D)$ denotes the irreducible complex characters of $D$.
More generally, for a normal subgroup $R$ and $\theta\in\Irr(R)$,
$\Irr(D\mid\theta)$ consists of the irreducible characters whose
restriction to $R$ contains $\theta$.

\begin{theorem}\label{thm:ordinary}
Let $D$ be a finite group, let $C\leq Z(D)$ be cyclic, and let
$\lambda\in\Irr(C)$ be faithful. Suppose that
$\Delta\leq\Aut(D)\times\Omega$ preserves $C$ and $\lambda$ under
the paired action, and acts transitively on $\Irr(D\mid\lambda)$.
Then $D$ is soluble.
\end{theorem}

Preserving $\lambda$ is a condition on each pair $(a,\sigma)$;
neither component is required to preserve it separately.
Navarro and Rizo proved the corresponding solubility result for an
automorphism action, in a more general normal-subgroup setting
\cite[Theorem~A]{NavarroRizo}. Allowing Galois operations requires an
additional argument, because they need not preserve irreducible
Brauer characters for an arbitrary finite group.

Theorem~\ref{thm:ordinary} is proved simultaneously with its
$2$-Brauer analogue for groups whose soluble-radical quotient is a
direct product of copies of $\GL_3(2)$. Canonical lifts give Galois
stability of Brauer characters in this class. A single minimal
counterexample, chosen across both character types, produces an
irreducible affine action. Its classification forces characteristic
two and a radical quotient $\GL_3(2)^k$. This is impossible in the
Brauer case, where the affine section has odd order. In the ordinary
case, restriction to odd-order elements and quotienting by a
nontrivial central $2$-subgroup produce a smaller Brauer counterexample.

The affine classification and Galois stability are established first,
before the shared minimal-counterexample argument.
The last section passes from characters to twisted group algebras.
In positive characteristic, a normal complement and a lift over
Witt vectors reduce the assertion to characteristic zero.
The proof uses the classification of finite simple groups through
the cited prime-power-index and Schur-multiplier results.

\subsection*{Conventions and standard ingredients}

All groups in the character-theoretic arguments are finite.
A group is \emph{perfect} if it equals its commutator subgroup.
For a finite group $D$, the notation $D^{(\infty)}$ means the final,
stable term of its derived series; it is perfect, and
$D/D^{(\infty)}$ is soluble.
We use $F(D)$ for the Fitting subgroup, the largest nilpotent normal
subgroup. A \emph{component} is a subnormal perfect subgroup whose
quotient by its center is nonabelian simple. The layer $E(D)$ is
the product of the components. The generalized Fitting subgroup is
$F^*(D)=F(D)E(D)$, and its centralizer is contained in it
\cite{IsaacsGroups}.
For a prime $p$, the symbols $O_p(D)$ and $O_{p'}(D)$ denote,
respectively, the largest normal $p$-subgroup and the largest
normal subgroup of order prime to $p$.

For a prime $p$, a $p$-regular element is an element whose order is
prime to $p$, and $\IBr_p(D)$ denotes the irreducible $p$-Brauer
characters. Fix compatible identifications of roots of unity of
order prime to $p$ with complex roots of unity. Brauer characters
are then complex-valued functions on the $p$-regular elements.
For an ordinary character $\chi$, its restriction to these elements
is denoted by $\chi^0$. Pointwise application of $\sigma\in\Omega$
to a Brauer character is always defined as a function; its
irreducibility will be proved whenever it is needed.
The notation $\IBr_p(D\mid\theta)$ has the same restriction meaning
as in the ordinary case. For either kind of character of a normal
subgroup, its inertia group is
$I_D(\theta)=\{g\in D:\theta^g=\theta\}$.

We use ordinary and modular Clifford theory and Gallagher's
correspondence in their usual forms
\cite[Chapters~6 and~11]{Isaacs}\cite[Theorems~4.6--4.8]{Vallejo}.
Thus restriction to a normal subgroup is homogeneous up to
conjugation, induction from the inertia group gives the Clifford
correspondence, and an extension of an invariant irreducible
character identifies the characters above it with the irreducible
characters of the quotient. These assertions will be applied only
after the relevant invariance and extension conditions have been
checked.

For a finite perfect group $Q$, its universal central extension
$U\longrightarrow Q$ is finite and has central kernel the Schur
multiplier of $Q$. The group $U$ is perfect and has trivial Schur
multiplier; this property is called \emph{superperfectness}.
Consequently, for the trivial actions on the coefficient groups,
\[
 H^2(U,\mathbb C^\times)=0,\qquad
 H^2(U,\overline{\mathbb F}_2^\times)=0.
\]
An invariant irreducible representation across a quotient $U$
therefore extends, and the extension is unique because $U$ has
no nontrivial linear characters. Automorphisms of $Q$ lift uniquely
to $U$; these facts give the equivariance in the descent arguments
below \cite{Brown}.

We work in the full group $\Omega$, rather than selecting lifts
between finite cyclotomic Galois groups. Every action on a particular
finite character set has finite image. Inner automorphisms can
always be adjoined to a paired acting group: they fix ordinary
and Brauer characters of the whole group and fix every central
character.


\section{An affine classification}
\label{sec:affine}

We first isolate the finite permutation-group argument used in both
character-theoretic reductions. For a finite-dimensional vector space $V$
over $\mathbb F_p$, write $\operatorname{AGL}(V)$ for its affine group and
$\ell:\operatorname{AGL}(V)\to\operatorname{GL}(V)$ for linear projection.
We identify $V$ with the subgroup of translations. The linear part of an
affine group $E$ is called irreducible when $V$ is an irreducible
$\mathbb F_p\ell(E)$-module.

\begin{lemma}\label{lemma:affine-radical}
Let $E\leq\operatorname{AGL}(V)$ be transitive, with irreducible linear part,
where $V\ne0$. Suppose that $1\ne J\triangleleft E$ is perfect and that
every point stabilizer in $J$ is soluble. Then
$\operatorname{Rad}(J)=1$, and linear projection is injective on $J$.
\end{lemma}

\begin{proof}
The image of $O_p(J)$ under linear projection is a normal $p$-subgroup of
$\ell(E)$. A $p$-group acting in characteristic $p$ has nonzero fixed
vectors. Its fixed space is $\ell(E)$-invariant, so irreducibility makes
this image trivial. Thus $O_p(J)$ consists of translations.

The translation subgroup $J\cap V$ is an $\ell(E)$-invariant subspace,
and hence equals $0$ or $V$. In the latter case $J=VJ_0$, so
$J/V\cong J_0$ is soluble. It is also perfect, as a quotient of $J$,
and therefore trivial. This would make $J=V$ abelian and perfect,
contrary to $J\ne1$. Consequently $J\cap V=1$ and $O_p(J)=1$.

The group $O_{p'}(J)$ fixes an affine point: average the orbit of any
point, using that its order is invertible in $\mathbb F_p$. Its fixed-point
set is $E$-invariant because $O_{p'}(J)$ is normal in $E$. Transitivity
makes this set all of $V$, and faithfulness gives $O_{p'}(J)=1$.
If the soluble radical of $J$ were nontrivial, its Fitting subgroup would
contain a nontrivial subgroup $O_r(J)$ for some prime $r$. This contradicts
$O_p(J)=O_{p'}(J)=1$.
\end{proof}

\begin{theorem}\label{thm:affine}
Under the hypotheses of Lemma~\ref{lemma:affine-radical}, suppose in addition
that $J$ is a direct product of nonabelian simple groups. Then there is an
integer $k\geq1$ such that
\[
 p=2,\qquad J\cong\operatorname{GL}_3(2)^k.
\]
\end{theorem}

\begin{proof}
Put $n=\dim_{\mathbb F_p}V$ and write
$J=S_1\times\cdots\times S_k$, where each $S_i$ is nonabelian simple.
Since $J$ is normal in the transitive group $E$, its orbits have equal
lengths dividing $p^n$. Each factor $S_i$ is normal in $J$, so its
orbits within a $J$-orbit also have equal lengths dividing that $p$-power. The stabilizer of any point in $S_i$ is soluble,
and is proper because $S_i$ is not soluble.

The classification of simple groups with a soluble subgroup of prime-power
index gives the following possibilities for a simple factor and its orbit
length:
\[
 (A_5,5),\quad
 (\operatorname{PSL}_3(2),7),\quad
 (\operatorname{PSL}_3(3),13),\quad
 (\operatorname{PSL}_2(q),q+1).
\]
In the last case $q+1$ is a prime power. The actions are the natural
point or hyperplane actions in the first three cases, and the projective-line
action in the last case; see \cite[Lemma~6 and Corollary~7]{BastosSchneider}.
In particular, for a fixed simple group and a fixed prime $p$, all the
orbits have the same length.

The fixed points of an affine transformation form either the empty set
or an affine subspace. A nonempty fixed-point set therefore has $p$-power
cardinality. This elementary observation excludes three entries at once.
A $3$-cycle in the degree-five action of $A_5$ fixes two points. An
involution in either degree-seven action of $\operatorname{PSL}_3(2)$
fixes three points. In either degree-thirteen action of
$\operatorname{PSL}_3(3)$, the involution represented by
$\operatorname{diag}(-1,-1,1)$ fixes five points: four on its two-dimensional
eigenspace and one on its complementary eigenline. On $V$ the corresponding
counts are, respectively, $2\cdot5^a$, $3\cdot7^a$ and $5\cdot13^a$
for some $a\geq0$, none of which is a power of the relevant characteristic.

It remains to consider $\operatorname{PSL}_2(q)$. If $q$ is even, a
nonidentity split-torus element fixes exactly two projective points.
Since the prime dividing $q+1$ is odd, its total fixed-point count on
$V$ is again impossible. Hence $q$ is odd; the same classification gives $q$ prime and
\[
 p=2,\qquad q+1=2^m.
\]
Nonabelian simplicity excludes $q=3$, so $m\geq3$.

Choose an element $x$ of order $q$ in $S_i$. It fixes one point in each
projective-line orbit, and therefore $2^{n-m}$ points on $V$. Translating
one fixed point to the origin shows that
\begin{equation}\label{eq:affine-rank}
 \operatorname{rank}(\ell(x)-1)=m.
\end{equation}
The multiplicative order of $2$ modulo $q=2^m-1$ is $m$: no smaller
positive exponent can give a multiple of $q$. Thus every nontrivial
irreducible $\mathbb F_2\langle x\rangle$-module has dimension $m$.
Equation~\eqref{eq:affine-rank} says that the nontrivial eigencharacters
of $x$ constitute a single Frobenius orbit of length $m$, each occurring
once. The normalizer of $\langle x\rangle$ in $S_i$ induces the squares
in $\mathbb F_q^{\times}$, a subgroup of order $(q-1)/2$, on these
eigencharacters. This subgroup acts freely on the nontrivial characters
and must preserve that one orbit. Hence
\[
 \frac{q-1}{2}\leq m.
\]
Together with $q=2^m-1$ and $m\geq3$, this forces $m=3$ and $q=7$.
Thus every $S_i$ is isomorphic to $\operatorname{GL}_3(2)$, as required.
\end{proof}



\section{Galois stability in characteristic two}
\label{sec:galois-stability}

Put $S=\operatorname{GL}_3(2)$, and consider the class
\[
\mathcal C=\bigl\{G\text{ finite}:G/\operatorname{Rad}(G)\cong S^k
                         \text{ for some }k\geq0\bigr\}.
\]
Here $\operatorname{Rad}(G)$ is the soluble radical. In particular,
$\mathcal C$ contains all soluble groups. We shall use this class exactly as
defined; knowing only the composition factors of a group would not suffice.

We need three facts about $S$. Every proper subgroup of $S$ is soluble;
its Schur multiplier has order two; and its irreducible $2$-Brauer characters
have degrees $1,3,3,8$. Each of these four Brauer characters is the restriction
of an ordinary irreducible character of the same degree. Since the ordinary
degrees are $1,3,3,6,7,8$, these four lifts form a Galois-stable set.
These facts are recorded in the character-table and subgroup data for $S$
\cite{GAP}. Consequently the full cyclotomic Galois group permutes
$\operatorname{IBr}_2(S)$ and, by taking tensor products, also
$\operatorname{IBr}_2(S^k)$.

We also use canonical lifts for soluble groups. If $R$ is soluble, restriction
to odd-order elements gives a bijection
\[
B_{2'}(R)\longrightarrow\operatorname{IBr}_2(R).
\]
The set on the left is invariant under group automorphisms and under all
cyclotomic Galois automorphisms \cite[Theorems~1.25 and~4.30]{Vallejo}.
Thus an invariant Brauer character has an invariant canonical ordinary lift,
and taking canonical lifts commutes with Galois conjugation.

\begin{lemma}\label{lemma:class-closure}
The class $\mathcal C$ is closed under subgroups, quotients, and finite
central extensions.
\end{lemma}

\begin{proof}
First let $H\leq S^k$. Project $H$ onto those coordinates on which its
projection is all of $S$. The image is a subdirect product of copies of a
nonabelian simple group, hence a product of diagonal copies of $S$.
This standard consequence of Goursat's lemma follows by adding one
coordinate at a time: the new coordinate either supplies an independent
factor or is identified with an existing simple quotient.
The kernel embeds in a product of proper subgroups of $S$, so it is soluble.
Since the image has trivial soluble radical, the kernel is precisely
$\operatorname{Rad}(H)$.

For a subgroup of an arbitrary $G\in\mathcal C$, apply this observation to
its image in $G/\operatorname{Rad}(G)$. The additional kernel lies in
$\operatorname{Rad}(G)$ and is soluble. Quotient closure follows because
every quotient of $S^k$ is a product of some of its simple factors.
Finally, in a central extension of $G$, the inverse image of
$\operatorname{Rad}(G)$ is soluble and its quotient is $S^k$; it is therefore
the soluble radical of the extension.
\end{proof}

Brauer characters will be regarded as complex-valued functions on the
odd-order elements, using a fixed identification of odd-order roots of
unity with their modular counterparts. For an arbitrary group, pointwise
Galois conjugation need not preserve irreducible Brauer characters. The next
result establishes the precise stability needed here.

\begin{theorem}\label{thm:galois-stability}
Let $G\in\mathcal C$. For every $\varphi\in\operatorname{IBr}_2(G)$ and
every cyclotomic Galois automorphism $\sigma$, the function
$\varphi^\sigma$ belongs to $\operatorname{IBr}_2(G)$.
\end{theorem}

\begin{proof}
We argue by induction on $|G|$. For soluble groups the assertion follows
from the canonical lifting bijection. Set $R=\operatorname{Rad}(G)$,
choose a constituent $\theta$ of $\varphi_R$, and let
$T=I_G(\theta)$.

If $T<G$, modular Clifford theory writes $\varphi=\psi^G$ for a unique
$\psi\in\operatorname{IBr}_2(T\mid\theta)$. The closure lemma and induction
give $\psi^\sigma\in\operatorname{IBr}_2(T)$. Moreover
\[
I_G(\theta^\sigma)=I_G(\theta)=T,
\]
because Galois conjugation commutes with group conjugation and is invertible.
The restriction of $\psi^\sigma$ to $R$ is a multiple of $\theta^\sigma$.
Clifford theory therefore makes $(\psi^\sigma)^G$ irreducible, and induction
of characters commutes with Galois conjugation. Hence
$\varphi^\sigma=(\psi^\sigma)^G$ is irreducible.

Suppose now that $\theta$ is $G$-invariant, and let
$\chi\in B_{2'}(R)$ be its canonical lift. This ordinary character is also
$G$-invariant. Let $U\to G/R$ be the universal central extension, with
kernel $M$, and form the fibre product
\[
\widehat G=G\times_{G/R}U.
\]
The quotient $\widehat G/R$ is $U$. Since $U$ is superperfect, the obstruction
in $H^2(U,\mathbb C^\times)$ vanishes and $\chi$ extends to an ordinary
irreducible character $\widehat\chi$ of $\widehat G$. The extension is
unique because $U$ has no nontrivial linear characters. Furthermore,
$M=\ker(\widehat G\to G)$ is a central $2$-group: the multiplier of $S^k$
is the product of the multipliers of its factors.

Choose an invariant lattice affording $\widehat\chi$ in a sufficiently large
$2$-modular system, and reduce it modulo its maximal ideal. Its restriction
to $R$ has Brauer character $\chi^0=\theta$, so that restriction, and
therefore the entire reduced module, is irreducible. The central group $M$
acts on the ordinary representation by scalars of $2$-power order. These
scalars reduce to one. The reduced module thus descends to an irreducible
Brauer character $\tau$ of $G$ extending $\theta$.

Apply the same reasoning to $\widehat\chi^\sigma$. Its restriction
$\chi^\sigma$ is the canonical lift of $\theta^\sigma$, so its reduction
again descends to an irreducible character of $G$. This character is
$\tau^\sigma$. Indeed, an odd-order element $g\in G$ has a unique
odd-order lift $\widetilde g$ in $\widehat G$: the inverse image of
$\langle g\rangle$ is abelian, and its odd Hall subgroup maps isomorphically
onto $\langle g\rangle$. Thus
\[
\tau(g)=\widehat\chi(\widetilde g),\qquad
\tau^\sigma(g)=\widehat\chi^\sigma(\widetilde g).
\]

Finally, the modular Gallagher correspondence gives
\[
\operatorname{IBr}_2(G\mid\theta)
 =\{\tau\beta:\beta\in\operatorname{IBr}_2(G/R)\},
\]
with quotient characters inflated to $G$. The Galois stability of
$\operatorname{IBr}_2(S^k)$ now shows that every
$(\tau\beta)^\sigma=\tau^\sigma\beta^\sigma$ is irreducible.
\end{proof}

The Galois operations in this proof are performed on ordinary characters
before reduction to characteristic two.

\subsection*{Restriction transfers transitivity}

\begin{lemma}\label{lemma:restriction-transitivity}
Let $D/\Rad(D)$ be a direct product of copies of $\GL_3(2)$.
Let $C\leq Z(D)$ be cyclic, let $\lambda\in\Irr(C)$ be faithful,
and suppose that a paired group $\Delta$ preserves $C$ and $\lambda$
and is transitive on $\Irr(D\mid\lambda)$. Then $\Delta$ is transitive on
$\IBr_2(D\mid\lambda|_{C_{2'}})$, where $C_{2'}$ is the odd-order
Hall subgroup of $C$.
\end{lemma}

\begin{proof}
Write $C=C_2\times C_{2'}$, where $C_2$ is the Sylow $2$-subgroup,
and write $\lambda=\lambda_2\times\lambda_{2'}$.
Theorem~\ref{thm:galois-stability} ensures that $\Delta$ permutes
\[
 X=\Irr(D\mid\lambda),\qquad
 Y=\IBr_2(D\mid\lambda_{2'}).
\]
Restriction to odd-order elements defines a linear map
\[
 \mathbb C X\longrightarrow\mathbb C Y.
\]
We verify surjectivity explicitly. The space $\mathbb C X$ consists
of the class functions $f$ on $D$ satisfying
$f(cg)=\lambda(c)f(g)$ for $c\in C$. The corresponding description
of $\mathbb C Y$ uses class functions on the odd-order elements
and the character $\lambda_{2'}$ of $C_{2'}$.
For such a target function $f$, use the commuting primary
decomposition $g=g_2g_{2'}$ and set
\[
 \widetilde f(g)=
 \begin{cases}
  \lambda_2(g_2)f(g_{2'}),&g_2\in C_2,\\
  0,&g_2\notin C_2.
 \end{cases}
\]
This is conjugation-invariant. Since elements of $C$ are central,
multiplication by $c=c_2c_{2'}$ replaces the primary factors by
$c_2g_2$ and $c_{2'}g_{2'}$. Thus
$\widetilde f(cg)=\lambda(c)\widetilde f(g)$.
Its restriction to the odd-order elements is $f$, proving
surjectivity.

The actions on the two vector spaces are the \emph{complex-linear}
extensions of their permutation actions on irreducible characters.
Restriction intertwines these actions. Indeed, the decomposition
numbers in
\[
 \chi^0=\sum_{\varphi\in Y}d_{\chi\varphi}\varphi
\]
are integers. Applying $\sigma$ to character values preserves
these coefficients and gives the corresponding equality for
$\chi^\sigma$; group automorphisms commute with restriction as
well. Dualizing this surjection embeds the space of invariant
functionals on $\mathbb C Y$ into that on $\mathbb C X$.
For a permutation basis, this dimension is the number of orbits.
Since $X$ is transitive and $Y$ is nonempty, $Y$ has exactly one orbit.
\end{proof}



\section{One argument for ordinary and Brauer characters}
\label{sec:character-transitivity}

We prove Theorem~\ref{thm:ordinary} together with the following
modular analogue.

\begin{theorem}\label{thm:modular-transitivity}
Let $G\in\mathcal C$, let $C\leq Z(G)$ be cyclic of odd order, and let
$\lambda\in\Irr(C)$ be faithful. Suppose a group of paired
automorphism--Galois operations preserves $C$, fixes $\lambda$, and is
transitive on $\IBr_2(G\mid\lambda)$. Then $G$ is soluble.
\end{theorem}

\begin{lemma}\label{lemma:character-cover}
Work either with ordinary characters of finite groups, or with irreducible
$2$-Brauer characters of groups in $\mathcal C$. Let $R\triangleleft N$
and let $\theta$ be an $N$-invariant irreducible character of $R$ of the
chosen kind. If $N/R$ is perfect, there exist a finite perfect group
$D_0$, a cyclic central subgroup $C_0$, and a faithful
$\lambda_0\in\Irr(C_0)$, such that $D_0/C_0\cong N/R$ and the
characters of $N$ above $\theta$ correspond bijectively to the characters
of $D_0$ above $\lambda_0$. The correspondence is equivariant for paired
operations preserving $R$ and fixing $\theta$. In the Brauer case,
$D_0\in\mathcal C$ and $C_0$ has odd order.
\end{lemma}

\begin{proof}
Let $U\twoheadrightarrow N/R$ be the universal central extension, with
kernel $M$, and form the pullback $\widehat N=N\times_{N/R}U$.
Thus $\widehat N/R\cong U$, while the kernel of
$\widehat N\twoheadrightarrow N$ is the central subgroup $M$.
The group $U$ is finite and superperfect. Hence the character-extension
obstruction vanishes, in $H^2(U,\mathbb C^\times)$ or
$H^2(U,\overline{\mathbb F}_2^\times)$ respectively. The resulting
extension $\widehat\theta$ is unique, because any two extensions differ
by a linear character of the perfect group $U$.

Let $\kappa$ be the scalar character of $M$ in this extension. In the
Brauer case it has odd order and is viewed as a complex linear character
trivial on the Sylow $2$-subgroup of $M$. The ordinary or modular Gallagher
correspondence gives the characters above $\theta$ in $\widehat N$ by
tensoring its extension with irreducible representations of $U$. Such a
tensor product descends to $N$ precisely when its scalar character on
$M$ is trivial. We may therefore take
\[
 D_0=U/\ker\kappa,\qquad C_0=M/\ker\kappa,
\]
with $\lambda_0$ induced by $\kappa^{-1}$. These have the required
properties; in the Brauer case, class membership follows from
Lemma~\ref{lemma:class-closure}.

An automorphism of $N$ preserving $R$ induces one of $N/R$, which
lifts uniquely to $U$ and gives a natural automorphism of the pullback.
In the Brauer case $\widehat N\in\mathcal C$,
so Theorem~\ref{thm:galois-stability} ensures that paired transforms of
$\widehat\theta$ remain irreducible. A pair fixing $\theta$ therefore
fixes its unique extension, and hence fixes $\kappa$ and preserves its
kernel. The tensor-product correspondence and its descent are equivariant.
\end{proof}

\begin{proof}[Proof of Theorems~\ref{thm:ordinary}
and~\ref{thm:modular-transitivity}]
Suppose that one of the theorems fails. Among all nonsoluble configurations
of either kind choose one minimizing
\[
 (|G:C|,|G|)
\]
lexicographically. Throughout this proof, ``characters'' means characters
of its chosen kind. Adjoin inner automorphisms to the paired group
$\Delta$, and let $\mathcal A$ be its projection to $\Aut(G)$.
Invariance below means invariance under $\mathcal A$.
In the Brauer case all subgroups, quotients, and central extensions used
below belong to $\mathcal C$, and their characters are Galois-stable, by
Lemma~\ref{lemma:class-closure} and Theorem~\ref{thm:galois-stability}.

\smallskip\noindent\emph{Minimality and restriction.}
For every invariant normal subgroup $L$, the paired action is transitive
on the characters of $L$ above $\lambda|_{C\cap L}$. Indeed, each such
representation $\rho$ extends to $CL$ by
$\widetilde\rho(cl)=\lambda(c)\rho(l)$, using the corresponding
odd-order modular roots of unity in the Brauer case. This extension
is covered by a character of $G$. Transitivity on the covering characters and Clifford theory
connect any two constituents on $L$, using an inner conjugation to align
the constituents.

Apply this observation to the last derived subgroup $G^{(\infty)}$.
Its index over $C\cap G^{(\infty)}$ is at most $|G:C|$, and its order
is smaller if it is proper. Minimality would make it, and hence $G$,
soluble. Thus $G$ is perfect.
In the Brauer case every irreducible character kills $O_2(G)$.
Since $C$ is odd, quotienting by a nontrivial $O_2(G)$ strictly decreases
the central index and preserves nonsolubility. Therefore $O_2(G)=1$,
and in particular $Z(G)$ is odd.

We claim that $C=Z(G)$. Choose a linear character $\nu$ of $Z(G)$ above
$\lambda$. Its paired stabilizer is transitive on the characters of
$G$ above $\nu$. Factoring by $\ker\nu$ gives a faithful cyclic central
character and central index $|G:Z(G)|$. If $Z(G)>C$, minimality makes
this quotient soluble, and its central kernel then makes $G$ soluble.
This quotient action is well defined: Galois conjugation preserves the
kernel of a linear character, so the paired stabilizer preserves
$\ker\nu$.

Every proper invariant normal subgroup $L$ is soluble. If $CL<G$,
restriction and minimality applied to $CL$ give this assertion. The
alternative $CL=G$ would make the perfect quotient $G/L$ abelian and
force $L=G$. Consequently $G/\Rad(G)$ is characteristically simple:
the inverse image of any proper characteristic subgroup is soluble,
while this quotient has trivial soluble radical. Thus
$G/\Rad(G)$ is a direct product of isomorphic nonabelian simple groups.

\smallskip\noindent\emph{Soluble inertia groups.}
Let $C<R<G$ be invariant and normal, and let $\theta$ be a character of
$R$ above $\lambda$. Put $T=I_G(\theta)$ and $N=RT^{(\infty)}$.
Clifford induction makes the paired stabilizer $\Delta_\theta$
transitive on the characters of $T$ above $\theta$: map their induced
characters in $G$, correct by an inner conjugation to fix $\theta$,
and use uniqueness of the Clifford correspondent. This stabilizer
preserves $N$ and contains the inner automorphisms of $T$. Restriction
and Clifford theory therefore give transitivity on the characters of
$N$ above $\theta$.

The quotient $N/R$ is perfect. Lemma~\ref{lemma:character-cover} produces
a configuration of the same kind whose central index is
\[
 |D_0:C_0|=|N:R|\leq |G:R|<|G:C|.
\]
Minimality makes $D_0$ soluble, and its perfectness gives $D_0=1$.
Thus $T^{(\infty)}\leq R$. Since $R$ is soluble, so is $T$.
In particular, no such $R$ has a unique character above $\lambda$,
because its inertia group would be $G$.

\smallskip\noindent\emph{Eliminating the layer.}
In the ordinary case we first show
\begin{equation}\label{eq:prime-support}
 \text{every prime dividing $|G|$ divides $|C|$.}
\end{equation}
Otherwise, for a prime $q\mid |G|$ with $q\nmid |C|$, choose a Sylow
$q$-subgroup $P$. The subgroup $CP=C\times P$ has linear character
$\rho=\lambda\times1_P$. Pairs normalizing $P$ fix $\rho$, and are
still transitive on $\Irr(G\mid\lambda)$: Sylow conjugacy adjusts any
pair by an inner automorphism without changing its action on characters
of $G$. If there are $r$ characters above $\lambda$, all have the same
degree $d$ and the same positive multiplicity $b$ in $\rho^G$.
The degree and degree-square identities give
\[
 |G:CP|=brd,\qquad |G:C|=rd^2.
\]
The first index is prime to $q$, whereas the second is divisible by $q$,
a contradiction.

The layer $E(G)$ is characteristic and perfect, so it is either trivial
or all of $G$. Suppose that $G=E(G)$. Then $G$ is a central product of
quasisimple groups, its radical is its center, and all the components
have the same simple central quotient.
In the ordinary case this simple group has a prime divisor absent from
its Schur multiplier \cite[Theorem~3.4]{NavarroRizo}. The centers of
the components are quotients of that multiplier and generate $Z(G)$.
This prime therefore divides $|G|$ but not $|C|$, contradicting
\eqref{eq:prime-support}.
In the Brauer case the simple quotient is $S=\GL_3(2)$, whose multiplier
has order two. Hence $Z(G)$ is a $2$-group. Since $O_2(G)=1$, we obtain
$G\cong S^k$ and $C=1$. Its trivial Brauer character is fixed by every
pair, whereas nontrivial irreducible Brauer characters exist. This again
contradicts transitivity. Thus $E(G)=1$ in both cases.

\smallskip\noindent\emph{An abelian normal section.}
The generalized Fitting subgroup is now $F=F(G)$, so $C_G(F)\leq F$.
It follows that $C<F<G$: equality $F=C$ would give $G\leq F$, and
$F=G$ would make $G$ soluble. In the Brauer case $F$ is odd because
$O_2(G)=1$.

In each Sylow subgroup of $F$ choose a characteristic Thompson critical
subgroup. Their product $T$ is characteristic, has class at most two,
and satisfies $C_F(T)=Z(T)$ \cite{Gorenstein}.
If $Z(T)\leq C$, set $R=CT$. Then $Z(R)=C$ and $C<R<G$.
Indeed, $T\leq C$ would give $F=C_F(T)=Z(T)\leq C$.
For $\chi\in\Irr(R\mid\lambda)$ and $x\notin C$, choose $y\in R$
with $[x,y]\ne1$. This commutator lies in $C$, so
\[
 \chi(x)=\chi(x^y)=\lambda([x,y])\chi(x).
\]
Faithfulness gives $\lambda([x,y])\ne1$, so $\chi(x)=0$.
Orthogonality gives $\chi(1)^2=|R:C|$, and the degree-square sum above
$\lambda$ then makes $\chi$ unique. In the Brauer case $R$ is odd, so
ordinary and Brauer characters coincide. In either case uniqueness
contradicts the inertia-group conclusion.

Therefore $CZ(T)>C$. Choose $A/C$ minimal nontrivial invariant in
$CZ(T)/C$. Then $A$ is abelian, $C<A<G$, and $A/C$ is an elementary
abelian $q$-group irreducible under $\mathcal A$. In the Brauer case
$q$ is odd.

\smallskip\noindent\emph{The affine contradiction.}
The set $\Irr(A\mid\lambda)$ is an affine space with translation space
$\operatorname{Hom}(A/C,\mu_q)$, where $\mu_q$ denotes the complex
$q$-th roots of unity. In the Brauer case these are also the Brauer
characters of $A$ above $\lambda$. Restriction gives a transitive paired
action with finite image $E$. Its linear parts are dual automorphisms
of $A/C$ multiplied by Galois scalars, so its linear part is irreducible.

Let $J$ be the image of the inner automorphisms of $G$ in $E$.
It is perfect and normal. It is nontrivial because the characters above
$\lambda$ separate $A$: their ratios separate $A/C$, and $\lambda$
is faithful on $C$. If $J=1$, every commutator $[a,g]$ would lie in all
their kernels, giving $A\leq Z(G)=C$.
The point stabilizers of $J$ are images of the soluble inertia groups
$I_G(\nu)$, with $\nu\in\Irr(A\mid\lambda)$.
Lemma~\ref{lemma:affine-radical} therefore gives $\Rad(J)=1$.
The kernel of $G\twoheadrightarrow J$ is a proper invariant normal
subgroup, hence soluble; since $J$ has trivial radical, that kernel is
exactly $\Rad(G)$. Thus $J$ is a direct product of nonabelian simple
groups, and Theorem~\ref{thm:affine} gives
\[
 q=2,\qquad G/\Rad(G)\cong J\cong\GL_3(2)^k\quad(k\geq1).
\]
This immediately contradicts the oddness of $q$ in the Brauer case.

It remains to consider the ordinary case. We have now proved
$G\in\mathcal C$, so Lemma~\ref{lemma:restriction-transitivity}
transfers the paired transitivity to
$\IBr_2(G\mid\lambda|_{C_{2'}})$.
Since $|G|$ is even, \eqref{eq:prime-support} shows that the Sylow
$2$-subgroup $C_2$ of $C$ is nontrivial. Every irreducible $2$-Brauer
character kills $C_2$, so the paired action descends to the modular
configuration
\[
 G/C_2,\qquad C/C_2,\qquad \overline\lambda,
\]
where $\overline\lambda$ is the faithful character induced by
$\lambda|_{C_{2'}}$. The action descends because $C_2$ is characteristic
in the preserved subgroup $C$. The quotient is still nonsoluble and belongs to
$\mathcal C$.

Its parameters satisfy
\[
 |G/C_2:C/C_2|=|G:C|,\qquad |G/C_2|<|G|.
\]
This contradicts our common minimal choice and proves both theorems.
\end{proof}


\section{Twisted group algebras over arbitrary fields}
\label{sec:fields}

\subsection{Characteristic zero}

Assume that $\operatorname{char}K=0$ and that $K^\alpha G$ is simple.
Let $k$ be the subfield generated over $\mathbb Q$ by the finitely
many coefficients of $\alpha$. The algebra $k^\alpha G$ is simple:
extending a nonzero proper ideal to $K$ would give a nonzero proper
ideal of $K^\alpha G$, by faithful flatness of field extensions.
The field $k$ is finitely generated over $\mathbb Q$, so we may
embed an algebraic closure $\overline k$ into $\mathbb C$.

Put $n=|G|$ and $b(g)=\prod_{x\in G}\alpha(g,x)$.
Multiplying the cocycle identity over $x\in G$, and using the
permutation $x\mapsto hx$, gives
\[
 \alpha(g,h)^n b(gh)=b(g)b(h).
\]
Choose $t(g)\in\overline k^\times$ with $t(g)^n=b(g)$ and $t(1)=1$.
Then
\[
 \beta(g,h)=\alpha(g,h)\frac{t(gh)}{t(g)t(h)}
\]
is a normalized cocycle with values in the group $\mu_n$ of
$n$-th roots of unity. Define the finite group
\[
 H=\mu_n\times_\beta G,\qquad
 (z,g)(w,h)=(zw\beta(g,h),gh).
\]
Its subgroup $\mu_n$ is central and $H/\mu_n\cong G$.
Inside the units of $\overline k^\alpha G$, the element $(z,g)$
is represented by $z u_g/t(g)$. Hence, with
$\lambda_0(z)=z$, the simple $\overline k^\alpha G$-modules
correspond to $\Irr(H\mid\lambda_0)$.

The algebra $k^\alpha G$ is finite-dimensional and simple.
Its center is a finite separable extension of $k$, and its
simple components after extension to $\overline k$ form one
orbit under $\Gal(\overline k/k)$. To describe the resulting
action on characters correctly, let
$\sigma\in\Gal(\overline k/k)$ and set
$\ell_\sigma(g)=\sigma(t(g))/t(g)\in\mu_n$.
Applying $\sigma$ to the embedded group of units gives the
automorphism
\[
 f_\sigma(z,g)
   =\bigl(\sigma(z)\ell_\sigma(g)^{-1},g\bigr).
\]
Thus the transitive action on $\Irr(H\mid\lambda_0)$ is
\[
 \chi\longmapsto\chi^\sigma\circ f_\sigma^{-1}.
\]
This formula follows by transporting a representation along the
semilinear action on $\overline k^\alpha G$.

Because these automorphisms arise from an actual Galois action
on the algebra, the map
\[
 \sigma\longmapsto
 \bigl(f_\sigma,\sigma|_{\mathbb Q^{\mathrm{ab}}}\bigr)
\]
is a homomorphism into $\Aut(H)\times\Omega$.
Its image preserves $\mu_n$ and fixes $\lambda_0$ as a pair:
on $\mu_n$, the first component acts by $\sigma$, which is
cancelled by the Galois operation on character values.
The transitivity just established therefore satisfies exactly
the hypotheses of Theorem~\ref{thm:ordinary}.
That theorem makes $H$ soluble, and its quotient $G$ is soluble.

\subsection{Positive characteristic}

We first record a lifting fact that makes the passage to
characteristic zero explicit.

\begin{lemma}\label{lemma:generic-simple}
Let $R$ be a discrete valuation ring with uniformizer $\pi$ and
fraction field $F$. Let $A$ be an $R$-algebra that is finite free
as an $R$-module. If $A/\pi A$ is simple, then $F\otimes_R A$ is
simple.
\end{lemma}

\begin{proof}
Suppose that $I$ is a nonzero proper two-sided ideal of
$F\otimes_R A$, and put $J=I\cap A$. Clearing denominators shows
that $FJ=I$. Thus $J$ has positive rank, and $A/J$ has positive
rank because $I$ is proper.
The submodule $J$ is saturated: if $\pi a\in J$ for $a\in A$,
then $a\in I$ and hence $a\in J$. In particular,
$J\cap\pi A=\pi J$ and $A/J$ is torsion-free.
Finitely generated torsion-free modules over a discrete valuation
ring are free. It follows that $J/\pi J$ is a nonzero proper
two-sided ideal of $A/\pi A$, a contradiction.
\end{proof}

Now assume $\operatorname{char}K=p>0$.
Passman's normal-complement and Hall-subalgebra results
\cite[Lemmas~2.6(ii) and~2.7]{Passman} give a normal subgroup
$N\triangleleft G$ such that $p\nmid|N|$, $G/N$ is a $p$-group,
and $K^\alpha N$ is simple. It is enough to prove that $N$ is
soluble.

Since $|N|$ is invertible in $K$, the algebra $K^\alpha N$ is
separable over $K$. One separability element is
\[
 \frac1{|N|}\sum_{g\in N}u_g\otimes u_g^{-1}.
\]
Indeed, its product is $1$, and multiplication by a basis unit
on either side agrees after reindexing the sum. The center of
$K^\alpha N$ is therefore a finite separable field extension
of $K$.

Let $k$ be the perfect closure of $K$. A finite separable
extension is linearly disjoint from a purely inseparable one;
thus the center remains a field after extension to $k$.
Since a central simple algebra stays central simple after
extending its center, $k^\alpha N$ remains simple.

The Witt ring $W(k)$ is a complete discrete valuation ring
of characteristic zero with residue field $k$ and uniformizer
$p$ \cite[Chapter~II]{Serre}. Its multiplicative Teichm\"uller map lifts $\alpha$
coefficientwise to a cocycle $\widetilde\alpha$ with unit
coefficients in $W(k)$. The algebra
$W(k)^{\widetilde\alpha}N$ is finite free, and its reduction
modulo $p$ is $k^\alpha N$.
Lemma~\ref{lemma:generic-simple} makes the twisted group algebra
of $N$ over the fraction field of $W(k)$ simple.
The characteristic-zero argument makes $N$ soluble.
Finally, $G/N$ is a $p$-group, so $G$ is soluble.
This proves Theorem~\ref{thm:main}.


\begin{thebibliography}{11}
\raggedright

\bibitem{BastosSchneider}
R.~Bastos and C.~Schneider,
\emph{Groups with a solvable subgroup of prime-power index},
Monatsh. Math. \textbf{194} (2021), 471--480.
\href{https://doi.org/10.1007/s00605-020-01486-5}{doi:10.1007/s00605-020-01486-5}.
See Lemma~6 and Corollary~7; the
\href{https://arxiv.org/abs/2006.16414}{author manuscript} is freely available.

\bibitem{Brown}
K.~S.~Brown,
\emph{Cohomology of Groups},
Graduate Texts in Mathematics, vol.~87,
Springer-Verlag, New York, 1982.
\href{https://doi.org/10.1007/978-1-4684-9327-6}{doi:10.1007/978-1-4684-9327-6}.

\bibitem{GAP}
T.~Breuer,
\emph{The GAP Character Table Library},
RWTH Aachen University, online database.
See the
\href{https://www.math.rwth-aachen.de/homes/Thomas.Breuer/ctbllib/ctbltoc/views/simplebyorder.html}{table of simple groups and Schur multipliers},
the \href{https://www.math.rwth-aachen.de/~Thomas.Breuer/ctbllib/ctbltoc/data/L3(2).html}{subgroup data for $\mathrm{L}_3(2)$},
and the
\href{https://www.math.rwth-aachen.de/~moc/decomposition/tex/L3(2)/L3(2)mod2.pdf}{characteristic-two decomposition matrix of $\mathrm{L}_3(2)$}.
Accessed 26 October 2026.

\bibitem{Gorenstein}
D.~Gorenstein,
\emph{On a theorem of Philip Hall},
Pacific J. Math. \textbf{19} (1966), no.~1, 77--80.
\href{https://doi.org/10.2140/pjm.1966.19.77}{doi:10.2140/pjm.1966.19.77}.

\bibitem{Isaacs}
I.~M.~Isaacs,
\emph{Character Theory of Finite Groups},
AMS Chelsea Publishing, vol.~359,
American Mathematical Society, Providence, RI, 2006.
Corrected reprint of the 1976 original.
\href{https://www.ams.org/bookpages/chel-359}{Publisher's information}.

\bibitem{IsaacsGroups}
I.~M.~Isaacs,
\emph{Finite Group Theory},
Graduate Studies in Mathematics, vol.~92,
American Mathematical Society, Providence, RI, 2008.
\href{https://doi.org/10.1090/gsm/092}{doi:10.1090/gsm/092}.

\bibitem{Kourovka}
E.~I.~Khukhro and V.~D.~Mazurov (eds.),
\emph{Unsolved Problems in Group Theory: The Kourovka Notebook},
21st ed., Novosibirsk, 2026,
Problem~18.78, p.~133.
\href{https://arxiv.org/abs/1401.0300v49}{arXiv:1401.0300v49},
20 October 2026.

\bibitem{NavarroRizo}
G.~Navarro and N.~Rizo,
\emph{Certain irreducible characters over a normal subgroup},
J. Group Theory \textbf{20} (2017), no.~4, 621--635.
\href{https://doi.org/10.1515/jgth-2016-0052}{doi:10.1515/jgth-2016-0052}.

\bibitem{Passman}
D.~S.~Passman,
\emph{Character theory and group rings},
in \emph{Character Theory of Finite Groups},
Contemporary Mathematics, vol.~524,
American Mathematical Society, Providence, RI, 2010, 139--148.
\href{https://doi.org/10.1090/conm/524/10352}{doi:10.1090/conm/524/10352}.
The \href{https://people.math.wisc.edu/~passman/charprob.pdf}{author manuscript}
contains the cited Lemmas~2.6 and~2.7.

\bibitem{Serre}
J.-P.~Serre,
\emph{Local Fields},
Graduate Texts in Mathematics, vol.~67,
Springer-Verlag, New York, 1979.
\href{https://doi.org/10.1007/978-1-4757-5673-9}{doi:10.1007/978-1-4757-5673-9}.

\bibitem{Vallejo}
C.~Vallejo Rodr\'iguez,
\emph{Characters, Correspondences and Fields of Values of Finite Groups},
doctoral dissertation, Universitat de Val\`encia, 2016.
\href{https://www.uv.es/cavaro3/download/Thesis.pdf}{Full text}.
See Theorems~1.25, 4.6--4.8 and~4.30.

\end{thebibliography}

\end{document}
